\documentclass[11pt,a4paper]{article}
\usepackage[T1]{fontenc}
\usepackage{lmodern,amsmath,amssymb,booktabs,microtype}
\usepackage[margin=25mm]{geometry}
\usepackage[colorlinks=true,urlcolor=blue,linkcolor=blue]{hyperref}
\newcommand{\tr}{\operatorname{tr}}
\newcommand{\Tr}{\operatorname{Tr}}
\newcommand{\E}{\mathbb E}
\newcommand{\Js}{\mathcal J_s}
\newcommand{\dd}{\mathrm d}
\title{The first derivative singlet in complex DSSYK\\Two moments in the double-scaled limit}
\author{Alan Day \and Chengheng\thanks{AI research assistant (OpenAI).}}
\date{Local research note --- 5 November 2026}
\begin{document}\maketitle
\begin{abstract}
We evaluate the ensemble-centered equal-time strength and the second
frequency moment of the first derivative singlet in complex DSSYK at
infinite temperature. Fermionic Wick contractions reduce finite-$N$
moments to coefficients of polynomials in at most six variables. Pairwise
set overlaps then yield an explicit double-scaled result. The centered
singlet has finite strength and finite initial evolution scale at fixed
positive double-scaling parameter. Its normalized initial curvature tends
to $16\Js^2$ in the subsequent small-parameter limit.
\end{abstract}

\section{A singlet with a finite dynamical scale}
The Susskind programme motivates seeking singlet observables that can
be matched to bulk propagation \cite{mss}. We use the same complex-SYK
model as in our preceding notes:
\begin{equation}
 H=\sum_{|I|=|J|=p}G_{IJ}c_I^\dagger c_J,\quad
 \E G_{IJ}G_{KL}=w\delta_{IL}\delta_{JK},\quad
 w=4^p\frac{\Js^2}{\lambda}\left(\frac{p!}{N^p}\right)^2.
\end{equation}
$G$ is Hermitian, $\{c_i,c_j^\dagger\}=\delta_{ij}$,
$\lambda=p^2/N$, $H=H_s$, and $\hbar=1$.
The state is the full $\beta=\mu=0$ trace, $\tr=2^{-N}\Tr$.
Define
\begin{equation}
 A=\sum_i[H,c_i]^\dagger[H,c_i],\qquad
 O=A-\E A,\qquad K=[H,A]=[H,O].
\end{equation}
Here $\E A=f(Q)$ is a deterministic function of particle number,
not the expectation value in one realization. Let
\begin{equation}
 V_0=\E\tr O^2,\qquad
 M_2=\E\tr K^\dagger K=\int\frac{\dd\omega}{2\pi}\omega^2S_O(\omega).
\end{equation}
For $N,p\to\infty$ with $p^2/N=\lambda>0$ fixed, our contraction
calculation gives, with $q=e^{-4\lambda}$,
\begin{align}
 \boxed{V_0^{\mathrm{DS}}=4\Js^4\left(\lambda^{-1}+1+q\right),}\label{v0}\\
 \boxed{M_2^{\mathrm{DS}}=\frac{8\Js^6}{\lambda^2}
       \left[1-q^2+4\lambda(q-q^2)\right].}\label{m2}
\end{align}
Both quantities are strictly positive for $\lambda>0$. The large mean
of $A$ has been removed, while its dynamics is unchanged.

\newpage
\section{What these moments tell us}
For each finite system, the normalized centered correlator obeys
\begin{equation}
 \frac{\E\tr O(t)O}{V_0}=1-\frac12\Omega^2 t^2+O_{N,p}(t^4),
 \qquad \Omega^2=\frac{M_2}{V_0},
\end{equation}
the limiting initial curvature is
\begin{equation}
 \Omega^2_{\mathrm{DS}}=
 \frac{2\Js^2}{\lambda}\,
 \frac{1-q^2+4\lambda(q-q^2)}{1+\lambda(1+q)}
 \ \longrightarrow\ 16\Js^2\quad(\lambda\to0^+).\label{curvature}
\end{equation}
Thus the candidate has a nontrivial initial evolution scale in string
units. This supplies a quantitative microscopic input for the
singlet-to-bulk programme. The next target is the frequency distribution
that realizes these moments, followed by its comparison with the bulk
Rindler propagator.

\paragraph{Exact finite-system information.}
The same calculation yields an exact coefficient formula for every
$N,p$. In particular, for $p=2$,
\begin{align}
 \frac{M_2(N,2)}{w^3}
 &=\frac{N^2(N-3)(N-2)(N-1)(N+1)}{512}\notag\\
 &\hspace{7mm}\times
   (N^4+4N^3-53N^2+120N-16).\label{p2}
\end{align}
It reproduces $120$ at $N=4$ and $2700$ at $N=5$ in $w=1$ units.
An independent full covariance-tensor sum gives the latter value.
For $p=1$, the dynamic moment vanishes identically at every $N$.

\paragraph{Approach to the limit.}
The table compares exact finite-$p$ coefficients with \eqref{v0} and
\eqref{m2}. These are ratios to the limiting expression, not fitted
extrapolations. At $\lambda=1/4$, $N=4p^2$; at $\lambda=1$, $N=p^2$.
\begin{center}
\begin{tabular}{rrrr}\toprule
$p$ & $V_0/V_0^{\rm DS}$, $\lambda=1/4$
    & $M_2/M_2^{\rm DS}$, $\lambda=1/4$
    & $M_2/M_2^{\rm DS}$, $\lambda=1$\\\midrule
2 & .940774 & .691904 & .218718\\
3 & .931858 & .803723 & .535200\\
4 & .937683 & .851019 & .595472\\
5 & .944487 & .878486 & .628144\\\bottomrule
\end{tabular}
\end{center}
The deviations remain substantial for these small $p$. The limiting
formulas follow from the overlap argument below, not from these data.

\section{Route through the literature}
We reuse fermionic Wick ordering \cite{ferialdi} and the embedded-ensemble
covariance structure \cite{kota,kv}. The known complex-SYK chord scale is
$q=e^{-4\lambda}$ \cite{bnr}. Particle-removal formulas in \cite{kv},
Section 7, organize neighboring charge sectors, but retain coupled
reduced matrix elements in the fourth mixed moment. For the present
sixth-order coupling average, we instead perform an explicit mode-wise
Wick reduction. This is an application of established contraction methods;
no claim of mathematical or physical priority is made.

\newpage
\section{Exact reduction to six variables}
Let $\mathcal D(X)=QX-\sum_i c_i^\dagger Xc_i$.
For a normal-ordered $k$-body monomial, $\mathcal D(X)=kX$.
When normal ordering a product of two $p$-body monomials $X,Y$, multiply
each contraction between the left annihilator block and the right creator
block by $1+h$. Call the resulting polynomial $R_h(X,Y)$. Then
\begin{equation}
 R_0(X,Y)=XY,\qquad
 \left.\partial_h R_h(H,H)\right|_0=(2p-\mathcal D)H^2=A.\label{mark}
\end{equation}
A term with $r$ such contractions has body rank $2p-r$, proving
\eqref{mark} directly.

Use six ordered Hamiltonian slots. For two disjoint adjacent marked
pairs $\mathcal M=((r,s),(t,u))$, set
\begin{equation}
 T_{ab}(h,g)=
 \begin{cases}\tfrac12,&a\le b,\\-\tfrac12,&a>b,\end{cases}
 \quad\text{with }T_{sr}\mapsto T_{sr}-h,\quad T_{ut}\mapsto T_{ut}-g.
\end{equation}
This is the signed thermal contraction matrix: a marked annihilator--creator
contraction changes from $1/2$ to $1/2+h$ or $1/2+g$.
For a coupling pairing $\pi$ (a fixed-point-free involution), define
\begin{equation}
 F_{\pi,\mathcal M}(z;h,g)=
 \sum_{S\subseteq\{1,\ldots,6\}}
   \det[T_{a,\pi(b)}]_{a,b\in S}\prod_{a\in S}z_a
 =\det(I+ZT P_\pi),\label{poly}
\end{equation}
where $Z=\operatorname{diag}(z_1,\ldots,z_6)$ and
$(TP_\pi)_{ab}=T_{a,\pi(b)}$. The empty determinant is one.
Let $[z^p]$ mean the coefficient of $z_1^p\cdots z_6^p$.
The contribution of a fixed paired word is
\begin{equation}
 w^3(-1)^{3p}[z^p]F_{\pi,\mathcal M}(z;h,g)^N.\label{finiteword}
\end{equation}
All overlaps $I\cap J$, including repeated mode labels in a vertex,
are included. No Fock-space cutoff or random-sample average is used.

Write $K^\dagger K=AHHA-AHAH-HAHA+HAAH$ before interpreting
each $A$ as a marked pair. The four choices are
\begin{equation}
 \begin{array}{c|c|c}
 &\mathcal M&\epsilon_{\mathcal M}\\\hline
 AHHA&((1,2),(5,6))&+1\\
 AHAH&((1,2),(4,5))&-1\\
 HAHA&((2,3),(5,6))&-1\\
 HAAH&((2,3),(4,5))&+1
 \end{array}
\end{equation}
Only pairings connecting the first three slots to the last three are
needed: the preceding note proves that $K$ has no linear Gaussian
component. Denote these six pairings by $\mathcal P_{\times}$.
\begin{equation}
 \boxed{\frac{M_2}{w^3}=(-1)^{3p}[z^p h g]
   \sum_{\pi\in\mathcal P_{\times}}
   \sum_{\mathcal M}\epsilon_{\mathcal M}F_{\pi,\mathcal M}^N.}\label{exact}
\end{equation}

\newpage
\appendix
\section{Signs, finite coefficients and independent checks}
For a pairing $\pi$, choose six independent creator sets $I_a$, each
of size $p$; the annihilator set at slot $b$ is $I_{\pi(b)}$.
Ordinary thermal Wick contraction is a determinant with $6p$ rows and
columns. Permute the annihilator blocks to the order matching the creator
blocks. The permutation has sign
$(-1)^{p^2(6/2)}=(-1)^{3p}$. Grouping both rows and columns by mode
then gives identical further permutation signs, whose product is one.
The determinant factors into the local minors in \eqref{poly}.
Summing the possible membership subsets of each of the $N$ modes gives
$F^N$, and $[z^p]$ enforces every set size. This proves \eqref{finiteword}.

For integer computation replace $T$ by $2T$, obtaining $G$ with integer
coefficients. The denominator of the extracted term is always $2^{6p}$.
Since $G-1$ has no constant term,
\begin{equation}
 \frac{M_2(N,p)}{w^3}=\frac{1}{2^{6p}}
             \sum_{j=0}^{6p}\binom Nj c_j(p),\label{binomial}
\end{equation}
where the signed $c_j$ are coefficients of $(G-1)^j$, summed over
the 24 marked configurations. For a fixed $p$, these coefficients are
computed once and then evaluated at any $N$ using exact integers.

An independent formula keeps all 15 Gaussian pairings in
$2\tr(H^2A^2-HAHA)$. Its coefficient arrays agree exactly with
\eqref{binomial} for $p=1,2,3$. The $p=1$ array is zero. Arrays through
$p=5$ vanish for all $N\le p+1$, as required by the conserved controls.
For $N=5,p=2$, direct Gaussian tensor contraction gives $2700$, agreeing
with \eqref{p2}.

The local marked-word formula itself was checked before ensemble
summation: 720 selected words at $(N,p)=(4,1),(5,2),(6,3)$ compare all
four coefficients $1,h,g,hg$ with direct CAR matrices. All 2880 rational
comparisons agree, including nonzero negative words. The selected index
sets and seed are retained; these are algebra tests, not disorder samples
used to estimate a moment.

\section{The double-scaled overlap limit}
For a cross pairing, the singleton coefficients of $F$ are
$a_i=+1/2$ for $i=1,2,3$ and $a_i=-1/2$ for $i=4,5,6$.
None is affected by our marks. Put
\begin{equation}
 W_S=\frac{\det[T_{a,\pi(b)}]_{a,b\in S}}{\prod_{a\in S}a_a},
 \qquad W_\varnothing=W_{\{i\}}=1.
\end{equation}
The coefficient in \eqref{finiteword}, divided by
$\binom Np^6\prod_i a_i^p$, is exactly
$\E_{I_1,\ldots,I_6}\prod_{\ell=1}^N W_{S_\ell}$ for six independent
uniform $p$-subsets of $\{1,\ldots,N\}$.

\newpage
Triple membership has probability at most
$\binom63p^3/N^2=o(1)$. The 15 pairwise overlap counts converge jointly
to independent Poisson variables of mean $\lambda$. One direct proof
uses joint factorial moments: assignments with distinct sites give
$(N)_r\prod_i(p)_{d_i}/(N)_{d_i}\to\lambda^r$, where
$\sum_i d_i=2r$; shared sites have higher membership and vanish.
This argument concerns a fixed number of sets, not a number growing
with $N$.

For completeness, convergence also holds for the weighted expectation,
not just in distribution. On a bounded complex neighborhood of $h=g=0$,
the finitely many $W_S$ obey $|W_S|\le R^{\binom{|S|}{2}}$ for some
$R\ge1$. If $Y=\sum_{i<j}|I_i\cap I_j|$, then
$|\prod W_S|\le R^Y$. Hypergeometric factorial moments and H\"older's
inequality give a uniform bound on $\E e^{cY}$ for any fixed $c$ when
$p^2/N$ remains in a compact interval. For one overlap $X$, for instance,
\begin{equation}
 \E e^{cX}\le
 \exp\!\left[\frac{p^2}{N-p+1}(e^c-1)\right],\qquad c\ge0.
\end{equation}
Cauchy--Schwarz bounds the contribution of triple membership by a
constant times $\sqrt{\binom63p^3/N^2}$, which tends to zero. The same
exponential-moment bound gives uniform integrability of the pair weights.
Consequently the convergence is locally uniform in complex $h,g$, and
Cauchy's formula permits the two marked derivatives to pass to the limit.

Write $b_{ij}$ for the coefficient of $z_i z_j$ in $F$. We obtain
\begin{equation}
 F^N\big|_{[z^p]}
 \sim\frac{N^{6p}}{(p!)^6}\prod_i a_i^p\,
 \exp\!\left\{\lambda\left[-3+
           \sum_{i<j}\left(\frac{b_{ij}}{a_i a_j}-1\right)\right]\right\},
\end{equation}
where this denotes convergence of the ratio by the nonzero displayed
prefactor, uniformly in the marking variables. With zero marks the
exponent is $-4\lambda\operatorname{cr}(\pi)$, reproducing the known
complex-SYK chord weight \cite{bnr}. In general it is
\begin{equation}
 -4\lambda\operatorname{cr}(\pi)+\lambda(\alpha h+\beta g+\delta hg).
\end{equation}
After the two derivatives, each graph contributes
\begin{equation}
 \frac{\Js^6}{\lambda^3}
 q^{\operatorname{cr}(\pi)}(\lambda\delta+\lambda^2\alpha\beta).
\end{equation}
In this normalization the factors $(-1)^{3p}$ and $\prod_i a_i^p$
cancel their signs, and $w^3$ cancels the factorial and $N$ powers.

\newpage
\section{The small graph sum}
The following table gives the coefficients after summing the four
commutator words, for each pairing. A tuple $(a,b,c)$ means that slots
$1,2,3$ are paired respectively with $a,b,c$.
\begin{center}
\begin{tabular}{rrrr}\toprule
Pairing & crossings & $\sum\epsilon\delta$
                         & $\sum\epsilon\alpha\beta$\\\midrule
$(4,5,6)$ & 3 & 0 & 0\\
$(4,6,5)$ & 2 & $-4$ & $-16$\\
$(5,4,6)$ & 2 & $-4$ & $-16$\\
$(5,6,4)$ & 1 & 0 & 16\\
$(6,4,5)$ & 1 & 0 & 16\\
$(6,5,4)$ & 0 & 8 & 0\\\bottomrule
\end{tabular}
\end{center}
Their sum is
\begin{equation}
 \frac{\lambda^3 M_2^{\mathrm{DS}}}{\Js^6}
 =8\lambda(1-q^2)+32\lambda^2(q-q^2),
\end{equation}
which proves \eqref{m2}.

The centered equal-time strength requires only four slots, with marks
$((1,2),(3,4))$. Subtracting $\E A$ removes the Gaussian pairing within
each marked pair. Of the two remaining pairings, the crossing one has
$(\alpha,\beta,\delta)=(-2,-2,0)$; the nested one has $(2,2,4)$.
The identical four-set overlap argument gives
\begin{equation}
 \frac{\lambda^2 V_0^{\mathrm{DS}}}{\Js^4}
 =4\lambda^2q+4\lambda+4\lambda^2,
\end{equation}
proving \eqref{v0}. Its finite coefficient calculation reproduces
$V_0(N,1)/w^2=N(N+1)^2/4$ and $V_0(4,2)/w^2=483/4$.

\section{Interpretation and limits}
These are model calculations for the full complex Gaussian ensemble,
with standard CAR and all overlapping index sets retained. The order
of limits in \eqref{curvature} is double scaling at fixed positive
$\lambda$, followed by $\lambda\to0^+$. Keeping $p$ fixed instead gives
a different limit; the exactly conserved $p=1$ model is a simple example.
No uniform error estimate for a simultaneous $\lambda(N)\to0$ limit
is asserted.

The Gaussian-centered $O$ can still overlap with sample-dependent
conserved operators. Thus $\Omega$ is the initial curvature scale of
this specified correlator, not an isolated particle mass, a late-time
decay rate, or a gap. The two moments do not determine the full spectral
density or establish convergence of the entire time correlator. Equation
\eqref{curvature} is a limit of finite-system curvatures; interchange of
the time and double-scaling limits requires further control.
In cosmic units $H_c=pH_s$ and $A_c=p^2A_s$,
$V_{0,c}=p^4V_{0,s}$ and $M_{2,c}=p^6M_{2,s}$.

\newpage
\section{Reproducibility and next calculation}
The integer coefficient engine and its drivers are in the code directory:
\begin{quote}\small\ttfamily
singlet\_coefficient.py\\
check\_singlet\_coefficient.py\\
run\_singlet\_double\_scale.py\\
run\_singlet\_centered\_moment.py
\end{quote}
The stage results retain exact arrays through $p=5$, factored polynomials,
all marked words, a separate full-tensor check, and the 24-graph table.
The numerical environment is Python 3.14.5, NumPy 2.4.6, SymPy 1.14.0.
Integer coefficient extraction has no statistical error or spectral
tail cutoff. The displayed finite-$p$ ratios are rounded only for reading.

An initial overlap audit omitted the $-1$ contribution of identically
zero minors when computing $\sum_{i<j}(b_{ij}/a_i a_j-1)$.
Its assertion failed and the original implementation was retained.
The repaired audit includes all 15 pairs and checks their known chord
exponents. The finite coefficient extraction and marked derivatives
were unaffected. The separate finishing script records the corrected
checks without restarting completed coefficient calculations.

The next calculation is a controlled construction of further moments
or the energy-resolved singlet kernel. The finite polynomial method
extends to longer words, while the marked overlap rules provide the
double-scaled simplification. We will first determine which part of
$V_0$ remains at zero frequency and compute the fourth frequency moment;
these are the next inputs needed to constrain the propagating spectrum.

\begin{thebibliography}{9}
\bibitem{mss} S. Miyashita, Y. Sekino and L. Susskind,
\emph{DSSYK at Infinite Temperature: The Flat-Space Limit and the
't Hooft Model},
\href{https://arxiv.org/html/2506.18054v3}{arXiv:2506.18054v3}.
\bibitem{ferialdi} L. Ferialdi and L. Di\'osi,
\emph{General Wick's Theorem for bosonic and fermionic operators},
Phys. Rev. A \textbf{104}, 052209 (2021),
\href{https://arxiv.org/html/2110.02920v2}{arXiv:2110.02920v2}.
\bibitem{kota} V. K. B. Kota,
\emph{SU(N) Wigner--Racah algebra for the matrix of second moments of
embedded Gaussian unitary ensemble of random matrices},
J. Math. Phys. \textbf{46}, 033514 (2005),
\href{https://arxiv.org/html/nucl-th/0410069v1}{arXiv:nucl-th/0410069v1}.
\bibitem{kv} V. K. B. Kota and M. Vyas,
\emph{Random Matrix Theory for Transition Strength Densities in Finite
Quantum Systems: Results from Embedded Unitary Ensembles},
\href{https://arxiv.org/html/1504.00759v1}{arXiv:1504.00759v1}.
\bibitem{bnr} M. Berkooz, V. Narovlansky and H. Raj,
\emph{Complex Sachdev-Ye-Kitaev model in the double scaling limit},
JHEP 02 (2021) 113,
\href{https://arxiv.org/html/2006.13983v2}{arXiv:2006.13983v2}.
\end{thebibliography}
\end{document}
